\section{Where this leaves a developer}
\label{sec:close}

\subsection{The through-line in one page}

The exposition began at a place every reader already knew and moved by
necessity rather than by exposition. It is worth collecting the steps, because
each was forced by the one before.

\begin{enumerate}[leftmargin=1.8em,itemsep=3pt]
\item The comprehension \texttt{SELECT-FROM-WHERE} is what data access looks
  like, in every language that has one. It has no \texttt{DO}.
\item LINQ exists to embed it in a host language so that a developer can finally
  express the \texttt{DO}. Rholang goes the rest of the way and builds the
  language around the comprehension:
  \texttt{for( ptrn <- chan where cond )P}, with \texttt{chan!(Q)} as its dual.
\item Making both the read and the write explicit turns a communication into a
  meeting, and the substitution the rule performs is the witness. Transactions
  become a property of the program text: a failed guard is not rolled back
  because nothing was consumed.
\item That leaves the pattern and the condition unaccounted for, which forces a
  data definition language. MeTTaIL's \texttt{language!} block supplies it in
  three rungs: grammar, equations, rewrites. The second rung holds all the
  invariants your type system cannot read; the third makes the data wiggle.
\item That structure subsumes MeTTa --- rules are rewrites, atom spaces are
  channels --- and adds the two things MeTTa lacks: an equational layer and a
  witnessed transaction.
\item The condition language is generated from the language definition rather
  than designed. Choosing the container for witnesses turns a dial between
  Boolean, linear and coalition sight, and the choice is an economic one about
  the deployment environment rather than a matter of rigour.
\item Grading the condition replaces the verdict with a value. At Booleans
  nothing changes; at the non-negative reals the program becomes a stochastic
  simulation of a population of itself; at the complex numbers, on a fenced
  fragment, a quantum one. Same text, two resolvers, no model drift.
\item In the graded setting the forager's guard is a learning rule, and the
  update is a correctly conditioned, unconditionally stable gradient step ---
  arrived at without importing anything.
\end{enumerate}

\subsection{What to do on Monday}

For a developer who wants to touch this rather than read about it, in order of
increasing commitment:

\begin{itemize}[leftmargin=1.6em,itemsep=3pt]
\item Run the two guard examples --- the receive guard and the \texttt{match}
  fall-through --- on the \texttt{cost-accounting-transpiler} branch of
  \texttt{f1r3node-rust}~\cite{f1r3node}. They are small, they compile, and they are covered by
  tests. Change the guard and watch what does and does not get consumed.
\item Run the cost-accounting demonstration in the same directory. It is a
  supply chain with conserved money and conserved inventory in which every
  operation is gated by a token, and it is the best available answer to what
  metering feels like when it is a program rather than a runtime flag.
\item Write a \texttt{language!} block for a data format you already own ---
  ideally one whose type has an invariant currently maintained by a comment. Put
  the comment in the \texttt{equations} block. That single exercise is the
  fastest way to feel the difference between rung one and rung two.
\item Then, if the shape of the thing appeals, add a \texttt{rewrites} block and
  discover that you have written a domain-specific language.
\end{itemize}

\subsection{Where to read next}

This note is deliberately the shallow end. The following are the deep ends, one
line each.

\begin{center}
\small
\begin{tabular}{@{}L{0.40\textwidth}L{0.54\textwidth}@{}}
\toprule
\textbf{Note} & \textbf{What it gives you that this does not} \\
\midrule
\emph{Graph-Structured Lambda Theories: A Reference for the Working
Developer}~\cite{GSLTintro} &
the same material with the theory first: what a language definition
\emph{is}, the category it lives in, and the two monads for cost and history \\[3pt]
\emph{Splitting the Mind of Zeus}~\cite{Zeus} &
why a logic has exactly the three components of \S\ref{sec:logic}, and where the
container dial comes from \\[3pt]
\emph{Observation Disciplines} (omnibus)~\cite{ObsDisc} &
the full range of the dial worked through six disciplines, with measurements,
and the argument that the disciplines compose but do not merge \\[3pt]
\emph{Graded Where-Clauses}~\cite{GradedWhere} &
the design and semantics of the graded slot, the resolution algebras, and the
result that weight tables are a normal form of graded clauses \\[3pt]
\emph{Races Decided by the Born Rule}~\cite{GWQuantum} &
the complex-valued case in detail, including the fence around it \\[3pt]
\emph{Graded Where-Clauses over the Reals}~\cite{GWPLN} &
the one-dimensional case: the gradient, the metric, and the measured forager \\[3pt]
\emph{The Mortal Scientist}~\cite{ScientistNote} &
the framework the running example is a small instance of --- learning as
cost-accounted computation \\
\bottomrule
\end{tabular}
\end{center}

\subsection{A closing observation}

There is a temptation, reading a note like this, to treat the last two sections
as the interesting part and the first two as preliminaries. Your author would
put it the other way round.

The two widenings are widenings of one slot. That slot exists because the
comprehension has a \texttt{WHERE} clause, and it has a \texttt{WHERE} clause
because SQL had one in 1974~\cite{ChamberlinBoyce74}. Nothing was invented to accommodate generated
logics or graded values; a place for them was already present in the shape every
developer already uses, and had been sitting empty. What this note describes is
mostly the consequence of taking that shape seriously enough to build a language
around it rather than embedding it in one.

The \texttt{DO} is the part LINQ went looking for. The rest of it is what you
find when you keep going.
